\documentclass[10pt,a4paper]{amsart}
\usepackage{amsmath,amssymb,amsthm,graphicx}
\usepackage{verbatim}
\newenvironment{reference}{\comment}{\endcomment}
%\newenvironment{reference}{}{}
\newenvironment{slogan}{\comment}{\endcomment}
\newenvironment{history}{\comment}{\endcomment}

% For commutative diagrams we use Xy-pic
\usepackage[all]{xy}

% We use 2cell for 2-commutative diagrams.
\xyoption{2cell}
% Theorem environments.
%
\theoremstyle{plain}
\newtheorem{theorem}[subsection]{Theorem}
\newtheorem{proposition}[subsection]{Proposition}
\newtheorem{lemma}[subsection]{Lemma}

\theoremstyle{definition}
\newtheorem{definition}[subsection]{Definition}
\newtheorem{example}[subsection]{Example}
\newtheorem{exercise}[subsection]{Exercise}
\newtheorem{situation}[subsection]{Situation}

\theoremstyle{remark}
\newtheorem{remark}[subsection]{Remark}
\newtheorem{remarks}[subsection]{Remarks}

\numberwithin{equation}{subsection}

% Macros
%
\def\lim{\mathop{\mathrm{lim}}\nolimits}
\def\colim{\mathop{\mathrm{colim}}\nolimits}
\def\Spec{\mathop{\mathrm{Spec}}}
\def\Hom{\mathop{\mathrm{Hom}}\nolimits}
\def\Ext{\mathop{\mathrm{Ext}}\nolimits}
\def\SheafHom{\mathop{\mathcal{H}\!\mathit{om}}\nolimits}
\def\SheafExt{\mathop{\mathcal{E}\!\mathit{xt}}\nolimits}
\def\Sch{\mathit{Sch}}
\def\Mor{\mathop{\mathrm{Mor}}\nolimits}
\def\Ob{\mathop{\mathrm{Ob}}\nolimits}
\def\Sh{\mathop{\mathit{Sh}}\nolimits}
\def\NL{\mathop{N\!L}\nolimits}
\def\CH{\mathop{\mathrm{CH}}\nolimits}
\def\proetale{{pro\text{-}\acute{e}tale}}
\def\etale{{\acute{e}tale}}
\def\QCoh{\mathit{QCoh}}
\def\Ker{\mathop{\mathrm{Ker}}}
\def\Im{\mathop{\mathrm{Im}}}
\def\Coker{\mathop{\mathrm{Coker}}}
\def\Coim{\mathop{\mathrm{Coim}}}

% Boxtimes
%
\DeclareMathSymbol{\boxtimes}{\mathbin}{AMSa}{"02}

%
% Macros for moduli stacks/spaces
%
\def\QCohstack{\mathcal{QC}\!\mathit{oh}}
\def\Cohstack{\mathcal{C}\!\mathit{oh}}
\def\Spacesstack{\mathcal{S}\!\mathit{paces}}
\def\Quotfunctor{\mathrm{Quot}}
\def\Hilbfunctor{\mathrm{Hilb}}
\def\Curvesstack{\mathcal{C}\!\mathit{urves}}
\def\Polarizedstack{\mathcal{P}\!\mathit{olarized}}
\def\Complexesstack{\mathcal{C}\!\mathit{omplexes}}
% \Pic is the operator that assigns to X its picard group, usage \Pic(X)
% \Picardstack_{X/B} denotes the Picard stack of X over B
% \Picardfunctor_{X/B} denotes the Picard functor of X over B
\def\Pic{\mathop{\mathrm{Pic}}\nolimits}
\def\Picardstack{\mathcal{P}\!\mathit{ic}}
\def\Picardfunctor{\mathrm{Pic}}
\def\Deformationcategory{\mathcal{D}\!\mathit{ef}}
\usepackage{cleveref}
\crefname{originallemma}{original lemma}{original lemmas}
\Crefname{originallemma}{Original lemma}{Original lemmas}
\makeatletter
\newlabel{algebra-lemma-quasi-finite-over-polynomial-algebra}{{00QE}{source}{}{sourcepointer-1}{}}
\newlabel{algebra-lemma-quasi-finite-over-polynomial-algebra@cref}{{[originallemma][0][]00QE}{[1][0][]source}{}{}{}}
\makeatother
\crefname{originaltheorem}{original theorem}{original theorems}
\Crefname{originaltheorem}{Original theorem}{Original theorems}
\makeatletter
\newlabel{algebra-theorem-openness-flatness}{{00RC}{source}{}{sourcepointer-2}{}}
\newlabel{algebra-theorem-openness-flatness@cref}{{[originaltheorem][0][]00RC}{[1][0][]source}{}{}{}}
\makeatother
\crefname{originallemma}{original lemma}{original lemmas}
\Crefname{originallemma}{Original lemma}{Original lemmas}
\makeatletter
\newlabel{algebra-lemma-flat-finite-presentation-limit-flat}{{02JO}{source}{}{sourcepointer-3}{}}
\newlabel{algebra-lemma-flat-finite-presentation-limit-flat@cref}{{[originallemma][0][]02JO}{[1][0][]source}{}{}{}}
\makeatother
\crefname{originallemma}{original lemma}{original lemmas}
\Crefname{originallemma}{Original lemma}{Original lemmas}
\makeatletter
\newlabel{lemma-base-change-gorenstein}{{0C07}{source}{}{sourcepointer-4}{}}
\newlabel{lemma-base-change-gorenstein@cref}{{[originallemma][0][]0C07}{[1][0][]source}{}{}{}}
\makeatother
\crefname{originallemma}{original lemma}{original lemmas}
\Crefname{originallemma}{Original lemma}{Original lemmas}
\makeatletter
\newlabel{lemma-upper-shriek-composition}{{0ATX}{source}{}{sourcepointer-5}{}}
\newlabel{lemma-upper-shriek-composition@cref}{{[originallemma][0][]0ATX}{[1][0][]source}{}{}{}}
\makeatother
\crefname{originallemma}{original lemma}{original lemmas}
\Crefname{originallemma}{Original lemma}{Original lemmas}
\makeatletter
\newlabel{lemma-shriek-affine-line}{{0AA1}{source}{}{sourcepointer-6}{}}
\newlabel{lemma-shriek-affine-line@cref}{{[originallemma][0][]0AA1}{[1][0][]source}{}{}{}}
\makeatother
\crefname{originallemma}{original lemma}{original lemmas}
\Crefname{originallemma}{Original lemma}{Original lemmas}
\makeatletter
\newlabel{lemma-affine-flat-Noetherian-gorenstein}{{0C08}{source}{}{sourcepointer-7}{}}
\newlabel{lemma-affine-flat-Noetherian-gorenstein@cref}{{[originallemma][0][]0C08}{[1][0][]source}{}{}{}}
\makeatother
\title{Gorensteinness of a flat finite presentation morphism is characterised by the existence of a Gorenstein quasi-finite morphism to affine space}
\author{The Stacks Project Authors}
\date{stacks-project:duality.tex:local-snapshot-20261012; GNU FDL 1.2 or later}
\begin{document}
\maketitle
\noindent\textit{Editorial note.} All delivered mathematical text is the author's own, copied from the pinned Stacks Project chapter file duality.tex; no statement or proof step was rewritten, reordered or removed.\par\medskip
\noindent\textit{Editorial note.} Preamble provenance: the pinned chapter file opens with an input of the shared preamble and carries no preamble text of its own, so the project's single Stacks-wide preamble staged at sources/stacks-duality/preamble.tex is excerpted verbatim under the recipe fixed by LANE\_BRIEF 22.1, as three spans, lines 8-13 (verbatim package and the reference/slogan/history environments), lines 14-18 (Xy-pic setup) and lines 154-217 (the native theorem declarations and mathematical macros). Nothing was added, removed or redefined; the wrapper class and the amsmath/amssymb/amsthm/graphicx base packages are supplied by the assembler.\par\medskip
\noindent\textit{Editorial note.} The counted statement and proof contain no cross reference, no citation and no included file, so the delivered body is self contained.\par\medskip
\noindent\textit{Editorial note.} Canonical location of the counted proposition: Stacks tag 0C09 (sources/stacks-duality/pinned-original-source.tex line 6768).\par\medskip
\noindent\textit{Allowed standard prerequisites (editorial).}\begin{itemize}
\item Existence of a quasi-finite morphism to affine space over the base. Algebra, Lemma labelled algebra-lemma-quasi-finite-over-polynomial-algebra, cited at proof lines 6790-6794. Produces the pair (U, g) and sets the value of d.
\item Openness of flatness and spreading out of flat finite presentation data. Algebra, Theorem labelled algebra-theorem-openness-flatness, cited at proof lines 6816-6820, and Algebra, Lemma labelled algebra-lemma-flat-finite-presentation-limit-flat, cited at proof lines 6820-6824. Allows assuming the polynomial algebra map flat and reduces to a Noetherian base.
\item Base change of the Gorenstein locus for flat finite type morphisms. Same chapter, Lemma labelled lemma-base-change-gorenstein, cited at proof lines 6824-6828. Justifies the reduction to the Noetherian case.
\item The upper shriek of the affine space projection and the Noetherian Gorenstein criterion. Same chapter, Lemmas labelled lemma-upper-shriek-composition and lemma-shriek-affine-line, cited at proof lines 6834-6838, and Lemma labelled lemma-affine-flat-Noetherian-gorenstein, cited at proof lines 6838-6841. Supplies f\textasciicircum{}!O\_S = g\textasciicircum{}!O\_\{A\textasciicircum{}n\_S\}[n] and concludes the equivalence.
\end{itemize}\medskip
\ifdefined\hypertarget\hypertarget{sourcepointer-1}{}\fi
\noindent\textit{Source reference.} Original lemma 00QE: sources/stacks-chow/tags.txt, source lines 964--964. Cross-chapter standard prerequisite located in the project tag list staged in the local snapshot; invoked at proof line 6793 of the counted proof. It is retained only as a navigation pointer; every load-bearing step is inside the delivered span.\par\medskip
\ifdefined\hypertarget\hypertarget{sourcepointer-2}{}\fi
\noindent\textit{Source reference.} Original theorem 00RC: sources/stacks-chow/tags.txt, source lines 998--998. Cross-chapter standard prerequisite located in the project tag list staged in the local snapshot; invoked at proof line 6813 of the counted proof. It is retained only as a navigation pointer; every load-bearing step is inside the delivered span.\par\medskip
\ifdefined\hypertarget\hypertarget{sourcepointer-3}{}\fi
\noindent\textit{Source reference.} Original lemma 02JO: sources/stacks-chow/tags.txt, source lines 3343--3343. Cross-chapter standard prerequisite located in the project tag list staged in the local snapshot; invoked at proof line 6816 of the counted proof. It is retained only as a navigation pointer; every load-bearing step is inside the delivered span.\par\medskip
\ifdefined\hypertarget\hypertarget{sourcepointer-4}{}\fi
\noindent\textit{Source reference.} Original lemma 0C07: sources/stacks-duality/pinned-original-source.tex, source lines 6614--6614. Same-chapter pointer into the author environment lemma at source line 6614; invoked at proof line 6822 of the counted proof. It is retained only as a navigation pointer; every load-bearing step is inside the delivered span.\par\medskip
\ifdefined\hypertarget\hypertarget{sourcepointer-5}{}\fi
\noindent\textit{Source reference.} Original lemma 0ATX: sources/stacks-duality/pinned-original-source.tex, source lines 4161--4161. Same-chapter pointer into the author environment lemma at source line 4161; invoked at proof line 6837 of the counted proof. It is retained only as a navigation pointer; every load-bearing step is inside the delivered span.\par\medskip
\ifdefined\hypertarget\hypertarget{sourcepointer-6}{}\fi
\noindent\textit{Source reference.} Original lemma 0AA1: sources/stacks-duality/pinned-original-source.tex, source lines 4609--4609. Same-chapter pointer into the author environment lemma at source line 4609; invoked at proof line 6838 of the counted proof. It is retained only as a navigation pointer; every load-bearing step is inside the delivered span.\par\medskip
\ifdefined\hypertarget\hypertarget{sourcepointer-7}{}\fi
\noindent\textit{Source reference.} Original lemma 0C08: sources/stacks-duality/pinned-original-source.tex, source lines 6722--6722. Same-chapter pointer into the author environment lemma at source line 6722; invoked at proof line 6840 of the counted proof. It is retained only as a navigation pointer; every load-bearing step is inside the delivered span.\par\medskip
\begin{lemma}
\label{lemma-flat-finite-presentation-characterize-gorenstein}
Let $f : X \to S$ be a morphism of schemes which is flat and locally
of finite presentation. Let $x \in X$ with image $s \in S$.
Set $d = \dim_x(X_s)$. The following are equivalent
\begin{enumerate}
\item $f$ is Gorenstein at $x$,
\item there exists an open neighbourhood $U \subset X$ of $x$
and a locally quasi-finite morphism $U \to \mathbf{A}^d_S$ over $S$
which is Gorenstein at $x$,
\item there exists an open neighbourhood $U \subset X$ of $x$
and a locally quasi-finite Gorenstein morphism $U \to \mathbf{A}^d_S$ over $S$,
\item for any $S$-morphism $g : U \to \mathbf{A}^d_S$
of an open neighbourhood $U \subset X$ of $x$ we have:
$g$ is quasi-finite at $x$ $\Rightarrow$ $g$ is Gorenstein at $x$.
\end{enumerate}
In particular, the set of points where $f$ is Gorenstein is open in $X$.
\end{lemma}

\begin{proof}
Choose affine open $U = \Spec(A) \subset X$ with $x \in U$ and
$V = \Spec(R) \subset S$ with $f(U) \subset V$. Then $R \to A$
is a flat ring map of finite presentation. Let $\mathfrak p \subset A$
be the prime ideal corresponding to $x$. After replacing $A$ by a
principal localization we may assume there exists a quasi-finite map
$R[x_1, \ldots, x_d]  \to A$, see
Algebra, Lemma \ref{algebra-lemma-quasi-finite-over-polynomial-algebra}.
Thus there exists at least one pair $(U, g)$ consisting of an
open neighbourhood $U \subset X$ of $x$ and a locally\footnote{If $S$
is quasi-separated, then $g$ will be quasi-finite.} quasi-finite morphism
$g : U \to \mathbf{A}^d_S$.

\medskip\noindent
Having said this, the lemma translates into the following algebra
problem (translation omitted). Given $R \to A$ flat and of finite
presentation, a prime $\mathfrak p \subset A$ and
$\varphi : R[x_1, \ldots, x_d] \to A$ quasi-finite at $\mathfrak p$
the following are equivalent
\begin{enumerate}
\item[(a)] $\Spec(\varphi)$ is Gorenstein at $\mathfrak p$, and
\item[(b)] $\Spec(A) \to \Spec(R)$ is Gorenstein at $\mathfrak p$.
\item[(c)] $\Spec(A) \to \Spec(R)$ is Gorenstein in an open neighbourhood
of $\mathfrak p$.
\end{enumerate}
In each case $R[x_1, \ldots, x_n] \to A$ is flat at $\mathfrak p$
hence by openness of flatness
(Algebra, Theorem \ref{algebra-theorem-openness-flatness}),
we may assume $R[x_1, \ldots, x_n] \to A$
is flat (replace $A$ by a suitable principal localization).
By Algebra, Lemma \ref{algebra-lemma-flat-finite-presentation-limit-flat}
there exists $R_0 \subset R$ and $R_0[x_1, \ldots, x_n] \to A_0$
such that $R_0$ is of finite type over $\mathbf{Z}$ and
$R_0 \to A_0$ is of finite type and $R_0[x_1, \ldots, x_n] \to A_0$ is flat.
Note that the set of points where a flat finite type morphism
is Gorenstein commutes with base change by
Lemma \ref{lemma-base-change-gorenstein}.
In this way we reduce to the case where $R$ is Noetherian.

\medskip\noindent
Thus we may assume $X$ and $S$ affine and that
we have a factorization of $f$ of the form
$$
X \xrightarrow{g} \mathbf{A}^n_S \xrightarrow{p} S
$$
with $g$ flat and quasi-finite and $S$ Noetherian. Then $X$ and
$\mathbf{A}^n_S$ are separated over $S$ and we have
$$
f^!\mathcal{O}_S = g^!p^!\mathcal{O}_S = g^!\mathcal{O}_{\mathbf{A}^n_S}[n]
$$
by know properties of upper shriek functors
(Lemmas \ref{lemma-upper-shriek-composition} and
\ref{lemma-shriek-affine-line}).
Hence the equivalence of (a), (b), and (c) by
Lemma \ref{lemma-affine-flat-Noetherian-gorenstein}.
\end{proof}

\end{document}
